\documentclass[10pt, oneside]{apuntes}

\title{Definición y esquemas para el Simulador}
\author{Pablo Martín Yuste}

\begin{document}
\maketitle

\section{Definición del simulador}
El simulador consiste en el fuselaje de la cabina de un planeador \todo{Modelo del planeador del simu}, montada sobre una junta universal que permite los giros de cabeceo y alabeo, y dotado de dos motores eléctricos que mueven las articulaciones en el corte del fuselaje.
\par El software empleado es \href{https://www.condorsoaring.com/}{Condor 3}, que permite una extración de los datos del estado de vuelo del planeador mediante UDP.
\par Este proyecto consiste en poder leer los datos extraídos del Condor, obtener las fuerzas y momentos que se sentirían en cabina, y cambiar la actitud del simulador para reproducirlos de la forma más fiel posible.

\section{Sistemas del simulador}
El funcionaimento del simulador se explica a continuación:
\begin{figure}[ht]
	\centering
	\begin{tikzpicture}
		%Nodos del diagrama de bloques
		\node[block] (condor3) {Condor 3};
		\node[block, below left=1cm of condor3] (decoder) {Decoder};
		\node[block, below=1cm of decoder] (physics) {Modelo de\\física};
		%Flechas del diagrama de bloques
	\end{tikzpicture}
\end{figure}

\end{document}
